\section*{Hoofdstuk \ref{ch:b2:muscle-movement} --- Spiercontractie en motoriek}

\begin{solution}{exo:b2:muscle-movement:1}
De dunne filamenten glijden langs de dikke, getrokken door myosinekoppen;
geen van beide filamenten verkort. De I-band en de H-zone worden smaller en
de Z-schijven naderen elkaar; de A-band (de dikke filamenten) behoudt zijn
lengte.
\end{solution}

\begin{solution}{exo:b2:muscle-movement:2}
(1) ATP gebonden, kop los; ATP gehydrolyseerd, kop gespannen. (2) De kop
bindt actine. (3) Fosfaat vrijgegeven, krachtslag, ADP vrijgegeven.
(4) Een nieuw ATP bindt en de kop laat los. ATP is nodig om de kop
\emph{los te maken}: zonder ATP blijft elke kop gebonden en slaat de spier
vast --- rigor.
\end{solution}

\begin{solution}{exo:b2:muscle-movement:3}
Zenuwimpuls; afgifte van acetylcholine en de eindplaatpotentiaal
(\qty{0.6}{ms}); actiepotentiaal van de spier over de vezel en via de
T-tubuli naar binnen (\qty{1}{ms}); calcium vrijgegeven uit het
\omterm{def:b2:cell-differentiation:fibre}{sarcoplasmatisch reticulum} (\qty{1}{ms}); calcium bindt troponine,
tropomyosine verschuift, koppen hechten zich (enkele milliseconden); de kracht
stijgt over enkele tientallen milliseconden.
\end{solution}

\begin{solution}{exo:b2:muscle-movement:4}
\omterm{def:b2:muscle-movement:motorunit}{Motorische eenheid}: één motorisch \omterm{def:b2:neurons-synapses:neuron}{neuron} en alle vezels die het innerveert.
\omterm{def:b2:muscle-movement:motorunit}{Enkelvoudige contractie}: de korte samentrekking door één impuls. \omterm{def:b2:muscle-movement:motorunit}{Tetanus}: de
versmolten, aanhoudende samentrekking door impulsen die sneller aankomen dan
de \omterm{def:b2:muscle-movement:motorunit}{enkelvoudige contractie} afneemt. De kracht wordt gedoseerd met de
vuurfrequentie van elke eenheid en met het aantal gerekruteerde eenheden.
\end{solution}

\begin{solution}{exo:b2:muscle-movement:5}
Elk half \omterm{def:b2:cell-differentiation:fibre}{sarcomeer} glijdt \qty{0.2}{\micro m} in \qty{50}{ms}:
\qty{4}{\micro m/s}; dat zijn 20 slagen van \qty{10}{nm} in
\qty{50}{ms}, 400 per seconde als één kop de hele tijd gebonden was.
\end{solution}

\begin{solution}{exo:b2:muscle-movement:6}
\qty{2}{cm} in \qty{0.1}{s}: \qty{0.2}{m/s}, 2 lengtes per seconde; elk
\omterm{def:b2:cell-differentiation:fibre}{sarcomeer} verkort $2/\num{40000}$ cm $= \qty{0.5}{\micro m}$ in
\qty{0.1}{s}, \qty{5}{\micro m/s}.
\end{solution}

\begin{solution}{exo:b2:muscle-movement:7}
$F/F_0 = 0.3125/(v/v_{\max} + 0.25) - 0.25$: bij 1, 3 en 6 lengtes per
seconde ($0.1$, $0.3$, $0.6$ van $v_{\max}$): $0.64$, $0.32$, $0.12$.
Vermogen $Fv$ (in $F_0\times$ lengtes per seconde): $0.64$, $0.95$, $0.71$
--- het grootst rond 3 lengtes per seconde, ongeveer een derde van $v_{\max}$.
\end{solution}

\begin{solution}{exo:b2:muscle-movement:8}
$80\times 30 = \qty{2400}{N}$; aan de voet $2400\times 5/45 =
\qty{270}{N}$.
\end{solution}

\begin{solution}{exo:b2:muscle-movement:9}
Bij \qty{10}{Hz} komen de impulsen elke \qty{100}{ms}, nadat het calcium van
de vorige verdwenen is: afzonderlijke contracties, die elk vanuit een
gedeeltelijk ontspannen toestand beginnen, wat een golving geeft. Bij
\qty{50}{Hz} komen ze elke \qty{20}{ms}, voordat het calcium is weggepompt:
het calcium blijft hoog en de kracht versmelt. De \omterm{def:b2:muscle-movement:motorunit}{tetanus} overtreft de
\omterm{def:b2:muscle-movement:motorunit}{enkelvoudige contractie} omdat één calciumpuls te kort is voor de dwarsbruggen
om de serie-elastische elementen op te rekken en de volle kracht te
ontwikkelen; aanhoudend calcium maakt dat mogelijk.
\end{solution}

\begin{solution}{exo:b2:muscle-movement:10}
Arbeid \qty{20}{kJ}, ATP-energie \qty{80}{kJ}, $1.6$ mol ATP. Eerste
\qty{5}{s}: $0.8$ mol \omterm{prop:b2:muscle-movement:energy}{creatinefosfaat}. De resterende $0.8$ mol ATP
via glycolyse: $0.4$ mol glucose, $0.8$ mol lactaat --- in
\qty{40}{L}, \qty{20}{mmol/L}.
\end{solution}

\begin{solution}{exo:b2:muscle-movement:11}
Geblokkeerd afgiftekanaal: geen calcium, geen contractie --- verlamming.
Geblokkeerde pomp: het calcium blijft hoog, de spier kan niet ontspannen ---
een contractuur. Ongevoelig troponine: tropomyosine verschuift nooit, geen
contractie. Geen calcium in het \omterm{def:b2:blood-circulation:blood}{bloed}: skeletspier trekt normaal samen (haar
calcium is intern); het \omterm{def:b2:heart:anatomy}{hart}, dat extern calcium nodig heeft om zijn afgifte
in gang te zetten, stopt.
\end{solution}

\begin{solution}{exo:b2:muscle-movement:12}
De kop zet tot de helft van de vrije energie van een ATP om in arbeid; de
rest is warmte, en de verliezen van de pompen, van de dwarsbruggen die bij
hoge snelheid hechten zonder te trekken, en van de elastische elementen
brengen het geheel op een kwart --- het rendement van een goede benzinemotor.
Twee manieren om kracht te doseren maken zowel fijne sturing bij een lage
inspanning mogelijk (kleine eenheden die één voor één worden toegevoegd en
elk sneller of trager vuren) als een groot bereik (een factor duizend, van één
kleine eenheid op het ritme van \omterm{def:b2:muscle-movement:motorunit}{enkelvoudige contracties} tot alle eenheden in
\omterm{def:b2:muscle-movement:motorunit}{tetanus}); een enkele schakelaar zou één kracht geven of geen.
\end{solution}

\begin{solution}{pb:b2:muscle-movement:1}
\textbf{1.} $v = \sqrt{2\times 9.81\times 0.4} = \qty{2.8}{m/s}$;
$\tfrac12\times 70\times 2.8^{2} = \qty{275}{J}$.
\textbf{2.} Arbeid $= 275 + 70\times 9.81\times 0.4 = \qty{550}{J}$ over
\qty{0.4}{m}: \qty{1370}{N}, twee keer het gewicht.
\textbf{3.} $F_0 = 500\times 30 = \qty{15000}{N}$; aan de voet
\qty{1500}{N}.
\textbf{4.} \qty{4}{cm} in \qty{0.3}{s}: \qty{0.13}{m/s}, $1.1$ lengtes per
seconde, $0.14\,v_{\max}$.
\textbf{5.} $F/F_0 = 0.3125/(0.14 + 0.25) - 0.25 = 0.55$: \qty{8300}{N}
in de spier, \qty{830}{N} aan de voet --- ver onder de benodigde
\qty{1370}{N}. Verkorting van de spier alleen kan deze sprong niet maken.
De tegenbeweging levert de rest: terwijl de springer door de knieën zakt,
slaan de uitgerekte pezen elastische energie op, en de spier, die bij hoge
kracht vrijwel isometrisch samentrekt, laadt ze op; daarna veren de pezen
sneller terug dan welke vezel ook kan verkorten.
\textbf{6.} $550/0.3 = \qty{1.8}{kW}$; \qty{92}{W/kg} strekspier.
\textbf{7.} $4/\num{48000}$ cm $= \qty{0.83}{\micro m}$ per \omterm{def:b2:cell-differentiation:fibre}{sarcomeer},
\qty{0.42}{\micro m} per half \omterm{def:b2:cell-differentiation:fibre}{sarcomeer}: 42 slagen van \qty{10}{nm}.
\textbf{8.} \qty{0.42}{\micro m} in \qty{0.3}{s}: \qty{1.4}{\micro m/s};
140 slagen per seconde bij voortdurende binding.
\textbf{9.} $500\times 6\times 10^{10}\times \num{48000}\times 300 =
4.3\times 10^{20}$ koppen.
\textbf{10.} $0.55\times \num{15000}\times 0.04 = \qty{330}{J}$; per
slag $0.4\times 8\times 10^{-20} = \qty{3.2e-20}{J}$; $1.0\times
10^{22}$ slagen.
\textbf{11.} $1.0\times 10^{22}/4.3\times 10^{20} = 24$ slagen per
kop, van de 42 beschikbare: ongeveer \qty{60}{\%}. De reserve is
klein --- de spier werkt dicht bij haar grens, zoals vraag 5 al liet zien.
\textbf{12.} $1.0\times 10^{22}/6.0\times 10^{23} = 0.017$ mol,
\qty{8.7}{g}; de energie ervan is $0.017\times 50 = \qty{860}{J}$, waarvan
\qty{330}{J} arbeid werd: \qty{39}{\%}.
\textbf{13.} Sneller dan één impuls per \qty{30}{ms}: boven
\qty{33}{Hz}; ongeveer 10 impulsen in \qty{0.3}{s}.
\textbf{14.} $9.9\times 10^{-6}\times 5\times 10^{-10}\times 6\times
10^{23} = 3\times 10^{9}$ vrije ionen; $1.5\times 10^{9}$ ATP.
\textbf{15.} Vezels: $500/(\pi\times(3\times 10^{-3})^{2}) = 1.8\times
10^{7}$; ATP van de dwarsbruggen per vezel per impuls $1.0\times
10^{22}/(1.8\times 10^{7}\times 10) = 6\times 10^{13}$, veertigduizend
keer de schatting voor de pomp. Het vrije calcium is een kleine rest: het
reticulum geeft in de orde van \qty{0.1}{mmol/L} af, tien keer de
vrije stijging, en bijna alles daarvan wordt meteen door troponine en de
buffers gebonden --- en alles moet worden teruggepompt.
\textbf{16.} Voor een maximale sprong vrijwel elke eenheid binnen enkele
tientallen milliseconden, de kleine eerst en de grote als laatste; een
rustige stap rekruteert misschien een tiende tot een vijfde ervan.
\textbf{17.} Eenheden worden in volgorde van grootte gerekruteerd: bij een
lage inspanning voegt elke nieuwe eenheid een kleine toename toe, bij een
hoge inspanning zijn de resterende eenheden groot en voegt elk een grote stap
toe.
\textbf{18.} De eindplaatpotentiaal van elke vezel wordt gehalveerd; de
vezels waarvan de potentiaal onder de drempel zakt, vuren niet, zodat de
\omterm{def:b2:muscle-movement:motorunit}{enkelvoudige contractie} en de \omterm{def:b2:muscle-movement:motorunit}{tetanus} kleiner zijn, en de sprong tekortschiet.
\textbf{19.} $5\times 20 = \qty{0.1}{mol}$ ATP: genoeg voor zes sprongen.
\textbf{20.} $20\times 20 = \qty{0.4}{mol}$ \omterm{prop:b2:muscle-movement:energy}{creatinefosfaat}:
ongeveer 23 sprongen.
\textbf{21.} Tien sprongen gebruiken $0.17$ mol, wat het \omterm{prop:b2:muscle-movement:energy}{creatinefosfaat} nog
niet uitput; na een stuk of twintig neemt de glycolyse het over, en hopen
lactaat en protonen zich op.
\textbf{22.} \qty{1.8}{kW} arbeid bij \qty{40}{\%} is \qty{4.6}{kW} aan
ATP, $0.09$ mol per seconde, wat $0.015$ mol zuurstof per seconde vraagt ---
\qty{0.34}{L/s}, \qty{21}{L/min}, tien keer wat het \omterm{def:b2:blood-circulation:blood}{bloed} kan aanvoeren:
de sprong is noodgedwongen anaeroob, en de zuurstof komt achteraf.
\textbf{23.} $0.17$ mol ATP vraagt $0.029$ mol zuurstof,
\qty{0.64}{L}; over een minuut terugbetaald \qty{640}{mL/min}, twee en een
half keer het rustverbruik.
\textbf{24.} ATP wordt binnen milliseconden opnieuw gefosforyleerd vanuit
\omterm{prop:b2:muscle-movement:energy}{creatinefosfaat}, zodat de concentratie nauwelijks daalt; rigor vraagt ATP
dicht bij nul, wat een levende spier met enige reserve nooit bereikt.
\textbf{25.} 42 beschikbare slagen per half \omterm{def:b2:cell-differentiation:fibre}{sarcomeer}, 24 nodig;
$F/F_0 = 0.55$ bij de snelheid van de sprong; $0.017$ mol, \qty{8.7}{g}, ATP.
\end{solution}
